\section*{Bab \ref{ch:g1:materials-around-us} --- Terbuat dari Apa Benda-Benda?}

\begin{solution}{exo:g1:materials-around-us:1}
Kursi adalah \omterm{def:g1:materials-around-us:object}{benda}. Kayu adalah \omterm{def:g1:materials-around-us:material}{bahan}: kursi dapat dibuat dari kayu.
\end{solution}

\begin{solution}{exo:g1:materials-around-us:2}
Kaca jendela: kaca. Paku: logam. Kerikil: batu. Buku: kertas. Kaus
kaki: kain (wol atau katun).
\end{solution}

\begin{solution}{exo:g1:materials-around-us:3}
Kayu: pensil, pintu (atau kursi, meja). Plastik: mangkuk, ember (atau
mainan, botol).
\end{solution}

\begin{solution}{exo:g1:materials-around-us:4}
Garpu logam: hanya garpu itu yang tidak terbuat dari kayu atau plastik.
(Dua \omterm{def:g1:materials-around-us:object}{benda} dari kayu dan satu dari plastik.)
\end{solution}

\begin{solution}{exo:g1:materials-around-us:5}
Kaca: kaca itu bening.
\end{solution}

\begin{solution}{exo:g1:materials-around-us:6}
Paku baja. Magnet menarik besi dan baja, bukan kaca atau kayu.
\end{solution}

\begin{solution}{exo:g1:materials-around-us:7}
Dua \omterm{def:g1:materials-around-us:object}{benda} terapung (balok kayu dan tutup plastik); tiga tenggelam
(kerikil, kelereng, dan paku).
\end{solution}

\begin{solution}{exo:g1:materials-around-us:8}
Bilahnya dari logam, pegangannya dari plastik. Gunting itu terbuat dari
dua \omterm{def:g1:materials-around-us:material}{bahan}.
\end{solution}

\begin{solution}{exo:g1:materials-around-us:9}
Di lingkaran biru: kaleng itu terbuat dari logam. Kaleng itu tidak
berada di lingkaran jingga karena magnet tidak menariknya: kaleng itu
terbuat dari aluminium, bukan dari besi atau baja.
\end{solution}

\begin{solution}{exo:g1:materials-around-us:10}
Jendela harus memasukkan cahaya, dan kita harus dapat melihat menembus
jendela. Kaca bening; kayu tidak.
\end{solution}

\begin{solution}{exo:g1:materials-around-us:11}
``Ya'': paku, penjepit kertas, dan kunci, 3 \omterm{def:g1:materials-around-us:object}{benda}. ``Tidak'': kerikil,
sendok, kaleng, kelereng, dan kaus kaki, 5 \omterm{def:g1:materials-around-us:object}{benda}.
$3 + 5 = 8$: setiap \omterm{def:g1:materials-around-us:object}{benda} mendapat tempatnya.
\end{solution}
